\section{Introdução}
\label{sec:introducao}

A administração pública brasileira exerce papel central na prestação de serviços essenciais à população. É no âmbito municipal que se concretizam as políticas públicas cotidianas, abrangendo desde a atenção primária à saúde e a educação básica até a manutenção da malha viária urbana e os serviços de saneamento. A execução contínua dessas atribuições demanda gestão eficiente e mobilização de expressivos recursos públicos, viabilizados tanto por receitas tributárias próprias quanto por transferências constitucionais das esferas estadual e federal.

Para atender às demandas da coletividade, o poder público recorre rotineiramente ao mercado para aquisição de insumos médico-hospitalares, contratação de obras de engenharia civil, locação de frotas e prestação de serviços especializados. Dessa forma, consolida-se a interface entre o ente público contratante (prefeituras e câmaras municipais) e os fornecedores privados. Embora a legislação nacional — estabelecida pela Nova Lei de Licitações e Contratos Administrativos (Lei Federal nº 14.133/2021)~\cite{brasil2021licitacoes}, pela Lei de Defesa da Concorrência (Lei Federal nº 12.529/2011)~\cite{brasil2011concorrencia} e pela Lei Anticorrupção (Lei Federal nº 12.846/2013)~\cite{brasil2013anticorrupcao} — imponha rígidos princípios de publicidade, isonomia e ampla competitividade, arranjos societários, familiares ou cadastrais entre supostos rivais podem comprometer a efetiva disputa sem deixar vestígios evidentes na documentação habilitatória individual.

Essa dificuldade decorre do fato de que instrumentos convencionais de fiscalização operam sob perspectiva predominantemente tabular, avaliando atas de habilitação, julgamento e empenhos de maneira fragmentada. Sob essa visão linear, verificações isoladas têm baixa capacidade de capturar padrões relacionais indiretos, tais como sócios administradores comuns distribuídos entre CNPJs distintos, empresas coligadas por \textit{holdings} familiares, compartilhamento de endereço comercial ou telefone entre concorrentes formais e coincidências de sobrenome que induzem falsos positivos em análises meramente nominais. O sinal investigativo útil emerge, portanto, da interseção estrutural entre a rede de co-disputa licitatória, a composição societária real e a execução financeira efetivada no exercício.

Para superar essa limitação, a Teoria dos Grafos oferece uma alternativa robusta para análises relacionais. Ao modelar órgãos licitantes, empresas participantes e seus sócios reais como vértices ($V$) de um grafo heterogêneo multicamada e ao representar participações, co-disputas simultâneas, vínculos de propriedade e indícios cadastrais como arestas ponderadas ($E$), a abordagem em redes permite podar o ruído periférico em tempo linear e isolar subestruturas fechadas — em especial triângulos societários ($K_3$) — que formalizam a sobreposição de controle econômico entre competidores.

A motivação prática para este estudo nasce da lacuna de granularidade observada nos portais locais de transparência: enquanto o portal municipal de Divinópolis/MG disponibiliza apenas visões agregadas~\cite{portal_transparencia_betha}, os dados jurisdicionados do Sistema Informatizado de Contas dos Municípios (SICOM), mantidos em dados abertos pelo Tribunal de Contas do Estado de Minas Gerais (TCE-MG)~\cite{tce_mg_dados_abertos,tcemg2025in01}, reúnem tabelas relacionais detalhadas de todo o ciclo da despesa pública, mas omitem os nomes e CPFs dos sócios e trazem chaves estrangeiras ausentes entre licitações e empenhos. Assim, este trabalho tem como objetivo geral desenvolver e validar um \textit{pipeline} computacional modular em Teoria dos Grafos para triagem indiciária de vínculos societários ocultos e materialidade contratual nas licitações do município de Divinópolis (MG) no exercício consolidado de 2025. Em estrita consonância com a jurisprudência do Tribunal de Contas da União (TCU)~\cite{tcu2011acordao2341,tcu2024acordao1798}, enfatiza-se que os escores e subgrafos gerados não constituem prova automática de fraude ou ilícito, mas sim instrumentos objetivos de priorização para auditoria humana em contraditório.

De modo específico, a investigação cumpre quatro objetivos operacionais encadeados: (i) saneamento canônico de documentos fiscais e construção de um grafo multicamada base a partir das tabelas licitatórias do SICOM; (ii) enriquecimento cadastral via Quadro de Sócios e Administradores (QSA) da Receita Federal por meio da BrasilAPI~\cite{brasilapi2026cnpj}, com \textit{cache} local e pontuação ponderada restrita a pares co-disputantes; (iii) redução determinística em listas de adjacência por indução de arestas, decomposição $1$-core via busca em largura (BFS), componentes conexas por Union-Find e enumeração de triângulos fechados $K_3$; e (iv) cruzamento com homologações, empenhos líquidos e pagamentos de 2025, aplicando filtro anti-falso positivo onomástico e pseudonimização integral sob a Lei Geral de Proteção de Dados Pessoais (LGPD)~\cite{brasil2018lgpd}.

O restante deste artigo estrutura-se da seguinte forma. A Seção~\ref{sec:referencial} apresenta o referencial teórico, reunindo o paradigma  jurídico-regulatório brasileiro, os estudos de redes societárias e homofilia em licitações, o formalismo multicamada e a justificativa algorítmica adotada. A Seção~\ref{sec:metodologia} detalha a arquitetura de análises e tratamento de dados, a formulação do escore composto, a cadeia de redução na classe anonimização LGPD. A Seção~\ref{sec:resultados} analisa os achados empíricos. Por fim, a Seção~\ref{sec:conclusao} consolida as conclusões, registra as limitações estruturais da base pública e delineia os trabalhos futuros.
